%--------------------------------------------------------------------------------
%
%
%--------------------------------------------------------------------------------
\pagebreak
\unnumberedSection{BN - Create a simulator breakpoint}

%--------------------------------------------------------------------------------
\begin{iPageDescEntry}{Syntax}
\texttt{bn  <modNum>, <type>, <adr> [ , <len> ]  }
\end{iPageDescEntry}

%--------------------------------------------------------------------------------
\begin{iPageDescEntry}{Description}

The \texttt{BN} command creates a simulator breakpoint.

The \texttt{modNum} parameter specifies the module for which the
breakpoint is active. A module number of \texttt{-1} specifies all
modules.

The \texttt{type} parameter specifies the type of breakpoint. The
following types are supported:

\begin{smallGapItemize}
	\item \texttt{CODE  } - Break when execution reaches the specified address.
    \item \texttt{DATA  } - Break on any data access to the specified address or range.
    \item \texttt{DATA\_R} - Break on a read access to the specified address or range.
    \item \texttt{DATA\_W} - Break on a write access to the specified address or range.
\end{smallGapItemize}
 
The \texttt{adr} parameter specifies the breakpoint address. For data
breakpoints, the optional \texttt{len} parameter specifies the length
of the address range. Data ranges of data breakpoint cannot overlap. 
Code breakpoints do not support address ranges. 

If a breakpoint with the same type and address range already exists for
another module, no new breakpoint entry is created. Instead, the
specified module is added to the existing breakpoint's module mask.
Newly created breakpoints are enabled automatically.

\end{iPageDescEntry}

%--------------------------------------------------------------------------------
\begin{iPageDescEntry}{Example}

The following examples demonstrate the \texttt{BN} command. 

\begin{lstlisting}[style=iPageOpStyle]
(0) ->bn 4, code, 0x1000
 (1) ->bn 4, code, 0x2000
 (2) ->bn 4, data_r, 0x3000
 (3) ->bl
 Idx  Type   State   Adr                  Modules         
 0    CODE   E       0x00_0000_1000       0000_0000_0000_0010
 1    CODE   E       0x00_0000_2000       0000_0000_0000_0010
 2    DATA_R E       0x00_0000_3000:0004  0000_0000_0000_0010
 (4) ->
\end{lstlisting}
\end{iPageDescEntry}

%--------------------------------------------------------------------------------
\begin{iPageDescEntry}{Notes}

Data breakpoint ranges are defined by their starting address and
length. The address is aligned to the corresponding range boundary.

\end{iPageDescEntry}